\section*{Capítulo \ref{ch:b2:phanerozoic-earth} --- A Terra fanerozoica}

\begin{solution}{exo:b2:phanerozoic-earth:1}
(a) Paris (\ang{5.1}~N) e Nova York (\ang{3.6}~S); a Cidade do Cabo e Sydney estavam perto de \ang{60}~S.
(b) Cidade do Cabo: $61.9 - 33.9 = \ang{28.0}$, $28.0 \times 111.2 = \qty{3114}{km}$; Nova York:
$40.7 + 3.6 = \ang{44.3}$, \qty{4926}{km}. (c) $3114/300 = \qty{10.4}{km}$ por milhão de anos,
\qty{1.0}{cm} por ano; $4926/300 = \qty{16.4}{km}$ por milhão de anos, \qty{1.6}{cm} por ano.
\end{solution}

\begin{solution}{exo:b2:phanerozoic-earth:2}
(a) 72, 120.1, 50.5 e \qty{135.4}{Myr}; média $(444 - 66)/4 = \qty{94.5}{Myr}$. (b) A do
fim do Permiano (Sibéria, a \qty{1.9}{Myr}), a do fim do Triássico (CAMP, \qty{0.4}{Myr}) e a do
fim do Cretáceo (Decã, \qty{0.5}{Myr}). (c) Não: os intervalos vão de 50 a
\qty{135}{Myr}; uma média existe para qualquer conjunto de datas.
\end{solution}

\begin{solution}{exo:b2:phanerozoic-earth:3}
(a) $\qty{91}{m}/\qty{0.036}{mm.yr^{-1}} = \num{2.5e6}$ anos. (b) $3.7/0.036 = 103$ vezes. (c)
Longe de qualquer costa, em água limpa, pouca lama dos rios chegava ao fundo; a chuva de
carapaças de algas era o sedimento principal.
\end{solution}

\begin{solution}{exo:b2:phanerozoic-earth:4}
(a) Do arenito fluvial ao giz os depósitos se aprofundam para cima: uma transgressão. A
areia de praia acima do giz registra uma regressão. (b) Não: ela registra uma subida do nível relativo do mar,
que a subsidência do local, ou uma queda do aporte de sedimento, também poderiam produzir. (c)
Encontrar a mesma transgressão, na mesma época, em outros continentes, ou remover a
subsidência por \omterm{def:b2:basins-compared:backstripping}{backstripping}.
\end{solution}

\begin{solution}{exo:b2:phanerozoic-earth:5}
(a) $233/97 = \qty{2.4}{m}$ por milhão de anos, \qty{2.4}{\micro m} por ano. (b) $3.7/0.0024 =
1540$: a subida atual é umas 1500 vezes mais rápida. (c) A curva obtida por \omterm{def:b2:basins-compared:backstripping}{backstripping}: as perfurações mostraram
que as cartas de ciclos acertam a cronologia das quedas mas exageram seu tamanho, por um fator de dois ou
mais; Nova Jersey tem seu próprio viés (topografia dinâmica), daí a estimativa de $100 \pm
\qty{50}{m}$.
\end{solution}

\begin{solution}{exo:b2:phanerozoic-earth:6}
(a) $\num{e5} \times \num{e9} \times 2900 = \qty{2.9e17}{kg}$ de basalto, e $\num{2.9e17}
\times 0.0075 = \qty{2.2e15}{kg}$ de dióxido de carbono, $\num{2.2e15}/0.04401 = \qty{4.9e16}{mol}$.
(b) $\num{278e-6} \times (44.01/28.97) \times \num{5.1e18} = \qty{2.15e15}{kg}$. (c) Um só pulso
liberou mais ou menos o que o ar inteiro continha em 1750. (d) $\num{e17}/\num{2.15e15} = 46.5$: quase cinquenta.
\end{solution}

\begin{solution}{exo:b2:phanerozoic-earth:7}
(a) $2630 \times \frac{\pi}{6} \times (\num{e4})^3 = \qty{1.38e15}{kg}$. (b) $\tfrac12 \times
\num{1.38e15} \times (\num{11.2e3})^2 = \qty{8.7e22}{J}$; a \qty{20}{km.s^{-1}},
\qty{2.76e23}{J}. (c) Um ano de \qty{47}{TW} são $\num{47e12} \times \num{3.156e7} =
\qty{1.48e21}{J}$: 59 e 186 anos.
\end{solution}

\begin{solution}{exo:b2:phanerozoic-earth:8}
(a) Cor escura, laminação fina, ausência de bioturbação, muita matéria orgânica, pirita; muitas vezes um desvio dos
isótopos de carbono e sinais de aquecimento. (b) O evento do Toarciano com Karoo--Ferrar, o OAE~1a com
Ontong Java. (c) As datas têm incertezas de até milhões de anos, e existem muitos eventos e
províncias, de modo que coincidências fortuitas são prováveis; um vínculo causal exige que a província comece
pouco antes do evento e um traçador do vulcanismo, como o mercúrio, nos mesmos sedimentos.
\end{solution}

\begin{solution}{exo:b2:phanerozoic-earth:9}
(a) $60.9 + 5.1 = \ang{66.0}$, $66.0 \times 111.2 = \qty{7339}{km}$ em \qty{200}{Myr}:
\qty{36.7}{km} por milhão de anos, \qty{3.7}{cm} por ano. (b) Mais que o dobro dos
\qty{1.6}{cm} por ano desde \qty{300}{Ma}. (c) Só a variação de latitude é conhecida; a variação de
longitude, livre num referencial paleomagnético, e qualquer desvio aumentam a distância.
\end{solution}

\begin{solution}{exo:b2:phanerozoic-earth:10}
(a) A fração de fundo de idade $t$ é $p(t) = \frac2T(1 - t/T)$. Então $\int_0^T tp\,\mathrm{d}t =
\frac2T\bigl(\frac{T^2}2 - \frac{T^2}3\bigr) = \frac T3$ e $\int_0^T \sqrt t\,p\,\mathrm{d}t =
\frac2T\bigl(\frac23 T^{3/2} - \frac25 T^{3/2}\bigr) = \frac8{15}\sqrt T$. (b) $T = 3 \times 64.2 =
\qty{192.6}{Myr}$; $2500 + \frac8{15} \times 350 \times 13.878 = \qty{5091}{m}$. (c) $T/2 =
\qty{96.3}{Myr}$: \qty{4332}{m}; a queda é de \qty{759}{m}, \qty{2}{\percent} menos que
\qty{774}{m}. (d) Modelo uniforme: $\frac23\sqrt\tau = \frac23\sqrt{2\bar t} = 0.943\sqrt{\bar t}$;
triângulo: $\frac8{15}\sqrt{3\bar t} = 0.924\sqrt{\bar t}$. A resposta depende da idade média muito
mais do que da forma.
\end{solution}

\begin{solution}{exo:b2:phanerozoic-earth:11}
(a) Em $\Delta s$, só o fator $1 - f^{-1/2}$ depende de $f$; divida-se por seu valor em $f = 2$,
$1 - 2^{-1/2} = 0.2929$. (b) $f = (1 - 0.2929\,\Delta s/432)^{-2}$: 1.071 para \qty{50}{m}, 1.239
para \qty{150}{m}; de 7 a \qty{24}{\percent} mais rápida que hoje. (c) $432 \times (1 -
0.8^{-1/2})/0.2929 = \qty{-174}{m}$: o mar desceria \qty{174}{m}. (d) A distribuição das idades
só muda à medida que o fundo antigo é subductado e substituído, o que leva um tempo da ordem de $\tau$,
cerca de \qty{100}{Myr}; o mar acompanha ao longo de dezenas de milhões de anos.
\end{solution}

\begin{solution}{exo:b2:phanerozoic-earth:12}
(a) $2^{-7.4/2} = 0.077$: \qty{92}{\percent} desapareceram. (b) $t = 2\log_2(1/0.19) = \qty{4.8}{Myr}$.
(c) $4.8/0.060 = 80$ vezes mais rápida. (d) Dentro de um andar longo, a extinção de fundo sozinha elimina
a maioria das espécies, de modo que um pulso curto se soma a um total grande e é difícil de separar dele; a
incompletude do registro também espalha as últimas ocorrências para trás.
\end{solution}

\begin{solution}{pb:b2:phanerozoic-earth:1}
\textbf{1.} 2500, 4250, 5304 e \qty{6000}{m}.
\textbf{2.} $\tau = 2 \times 64.2 = \qty{128.4}{Myr}$; $2500 + 350\sqrt{128.4} = \qty{6466}{m}$.
\textbf{3.} $\num{293e6}/(65000 \times 128.4) = \qty{35.1}{km}$ por milhão de anos,
\qty{35.1}{mm} por ano: uma taxa lenta.
\textbf{4.} $\frac1\tau\int_0^\tau(2500 + 350\sqrt t)\,\mathrm{d}t = 2500 + \frac23 \times
350\sqrt\tau = \qty{5144}{m}$.
\textbf{5.} A lei da raiz quadrada é profunda demais para o fundo com mais de algumas dezenas de milhões de anos,
que se aplaina; as idades reais não estão uniformemente distribuídas, sendo o fundo jovem o mais abundante; e
o sedimento, os \omterm{def:b2:oceanic-islands:seamount}{montes submarinos} e os platôs elevam o fundo.
\textbf{6.} $\num{293e6} \times 5.144 = \qty{1.51e9}{km^3}$, \qty{13}{\percent} a mais que o
oceano inteiro: o modelo é profundo demais.
\textbf{7.} $v = \qty{17.6}{km}$ por milhão de anos; $17.6 \times 128.4 = \qty{2260}{km}$ de
cada lado.
\textbf{8.} $t = 1000/17.6 = \qty{56.8}{Myr}$; $2500 + 350\sqrt{56.8} = \qty{5138}{m}$.
\textbf{9.} Num flanco, com $x = vt$: $\int_0^{v\tau}\bigl(d(\tau) - d(x/v)\bigr)\mathrm{d}x
= vc\int_0^\tau(\sqrt\tau - \sqrt t)\,\mathrm{d}t = \frac13 vc\,\tau^{3/2}$; o dobro para os dois.
$\frac23 \times 17.6 \times 0.350 \times 128.4^{3/2} = \frac23 \times 17.6 \times 0.350 \times
1454.9 = \qty{5975}{km^2}$.
\textbf{10.} $65000 \times 5975 = \qty{3.88e8}{km^3}$.
\textbf{11.} $\num{3.88e8}/\num{1.33e9} = 0.29$.
\textbf{12.} Seu fundo é mais jovem em média, portanto mais quente, menos denso e mais alto; a faixa de fundo
jovem ao longo do eixo é mais larga.
\textbf{13.} $\tau = \qty{64.2}{Myr}$; $2500 + \frac23 \times 350\sqrt{64.2} = \qty{4370}{m}$.
\textbf{14.} $5144 - 4370 = \qty{774}{m}$.
\textbf{15.} $\num{293e6} \times 0.774 = \qty{2.27e8}{km^3}$.
\textbf{16.} $\num{2.27e8}/\num{362e6} = \qty{0.627}{km}$, \qty{627}{m}.
\textbf{17.} $\Delta s = 627 \times (3300 - 1025)/3300 = \qty{432}{m}$.
\textbf{18.} $f = \bigl(1 - 0.2929 \times 100/432\bigr)^{-2} = 1.15$: \qty{15}{\percent} mais rápida.
\textbf{19.} $0.149 \times 148 = \qty{22.1}{\mega km^2}$.
\textbf{20.} $48.1 + 0.32 \times 7.7 = \qty{50.6}{\percent}$, \qty{74.9}{\mega km^2}.
\textbf{21.} À medida que sobe, ele cobre mais área, de modo que o mesmo volume forma uma camada mais fina
(\cref{prop:b2:phanerozoic-earth:flooding}).
\textbf{22.} Áreas inundadas a 100, 200, 300 e \qty{400}{m}: 22.1, 41.7, 57.3 e
\qty{71.2}{\mega km^2}. Degraus: $(362 + 11.0) \times 0.1 = \qty{37.3e6}{km^3}$, depois 39.4, 41.2
e \qty{42.6e6}{km^3}; acumulados 37.3, 76.7, 117.9 e \qty{160.5e6}{km^3}. O volume a colocar
é \qty{156.4e6}{km^3}, entre 300 e \qty{400}{m}: $300 + (156.4 - 117.9)/42.6 \times 100
= \qty{390}{m}$.
\textbf{23.} $38.7 + 0.90 \times 9.4 = \qty{47.2}{\percent}$ da terra atual.
\textbf{24.} Dobrar a taxa de expansão elevaria o mar em cerca de \textbf{\qty{390}{m}} e
afogaria quase metade da terra atual.
\end{solution}
